% =========================================================
% Geometric Modeling --- Parametric Curves and Surfaces
% =========================================================

\section{Parametric Curves and Surfaces}

% ---------------------------------------------------------
% Toolkit
% ---------------------------------------------------------

% ---------------------------------------------------------
\subsection*{Parametric Curves}

\begin{propositionbox}{Definition (Parametric Curve)}
A \emph{parametric curve} in $\mathbb{R}^n$ is a vector-valued function
\[
  C\,:\,[a,b]\;\to\;\mathbb{R}^n, \qquad
  C(u) = \bigl(x_1(u),\,x_2(u),\,\ldots,\,x_n(u)\bigr),
\]
where each coordinate function $x_i(u)$ is a real-valued function
of the scalar parameter $u$.
The interval $[a,b]$ is typically normalised to $[0,1]$.
\end{propositionbox}

\begin{propositionbox}{Definition (Velocity, Regularity, Unit Tangent)}
The \emph{velocity vector} (or \emph{tangent vector}) at $u$ is the derivative
\[
  C'(u) = \bigl(x_1'(u),\,x_2'(u),\,\ldots,\,x_n'(u)\bigr).
\]
The curve is \emph{regular} at $u$ if $C'(u) \neq \mathbf{0}$.
At a regular point, the \emph{unit tangent vector} is
\[
  \hat{T}(u) = \frac{C'(u)}{\|C'(u)\|}.
\]
The scalar $\|C'(u)\|$ is the \emph{speed} of traversal.
\end{propositionbox}

\begin{remark*}[Properties of parametric curves]
Parametric curves carry structure that implicit curves do not:
\begin{itemize}
  \item \textbf{Direction.}
        The parameter increases from $a$ to $b$, giving the curve
        a natural orientation.
        Reversing the parametrisation reverses the direction.
  \item \textbf{Velocity and speed.}
        $C'(u)$ measures both the direction of motion and how fast
        the curve is being traversed.
        A curve can be \emph{reparametrised} (change of variable
        $u = \phi(t)$) without changing its shape;
        the velocity changes but the trace does not.
  \item \textbf{Endpoint interpolation.}
        The start and end points are $C(a)$ and $C(b)$, accessible directly.
  \item \textbf{Boundedness.}
        The curve is defined on a closed interval, so it represents
        a bounded segment, not an infinite line.
  \item \textbf{Cusps.}
        Where $C'(u) = \mathbf{0}$, the curve may have a cusp
        (velocity drops to zero, direction may reverse).
        These are degenerate points excluded by the regularity condition.
\end{itemize}
\end{remark*}

% ---------------------------------------------------------
\subsection*{Parametric Surfaces}

\begin{propositionbox}{Definition (Parametric Surface)}
A \emph{parametric surface} in $\mathbb{R}^3$ is a vector-valued function
of two parameters,
\[
  S\,:\,R\;\longrightarrow\;\mathbb{R}^3, \qquad
  S(u,v) = \bigl(x(u,v),\;y(u,v),\;z(u,v)\bigr),
\]
where $R\subseteq\mathbb{R}^2$ is the \emph{parameter domain}
(typically the unit square $[0,1]^2$).
\end{propositionbox}

% ---------------------------------------------------------
\subsection*{Partial Derivatives and Isoparametric Curves}

\begin{propositionbox}{Definition (Partial Derivative Vectors)}
The \emph{partial derivative vectors} of $S$ at $(u,v)$ are
\[
  S_u(u,v) = \frac{\partial S}{\partial u}
  = \left(\frac{\partial x}{\partial u},\;
          \frac{\partial y}{\partial u},\;
          \frac{\partial z}{\partial u}\right),
  \qquad
  S_v(u,v) = \frac{\partial S}{\partial v}
  = \left(\frac{\partial x}{\partial v},\;
          \frac{\partial y}{\partial v},\;
          \frac{\partial z}{\partial v}\right).
\]
$S_u$ is the tangent vector to the surface in the $u$-direction:
it is the velocity of the curve obtained by fixing $v$ and varying $u$.
$S_v$ is the corresponding tangent vector in the $v$-direction.
\end{propositionbox}

\begin{propositionbox}{Definition (Isoparametric Curves)}
Fixing one parameter on $S(u,v)$ produces a curve on the surface:
\begin{itemize}
  \item \emph{$u$-curve} at $v = v_0$:\quad
        $C_{v_0}(u) = S(u,\,v_0)$,\quad tangent $= S_u(u,v_0)$.
  \item \emph{$v$-curve} at $u = u_0$:\quad
        $C_{u_0}(v) = S(u_0,\,v)$,\quad tangent $= S_v(u_0,v)$.
\end{itemize}
These \emph{isoparametric curves} (isocurves) sweep the surface and
form its parameter grid.
The two families intersect at every surface point $S(u_0,v_0)$.
\end{propositionbox}

% ---------------------------------------------------------
\subsection*{Directional Derivative}

\begin{propositionbox}{Definition (Directional Derivative on a Surface)}
The \emph{directional derivative} of $S$ at $(u_0,v_0)$ in the
direction $\mathbf{d} = (d_u, d_v)$ (a unit vector in parameter space) is
\[
  D_{\mathbf{d}}\,S = d_u\,S_u(u_0,v_0) + d_v\,S_v(u_0,v_0).
\]
It measures the rate of change of position on the surface as the
parameter moves in direction $\mathbf{d}$ in the $uv$-plane.
Every tangent vector to the surface at $(u_0,v_0)$ is a directional
derivative for some choice of $\mathbf{d}$.
\end{propositionbox}

% ---------------------------------------------------------
\subsection*{Unit Normal and Tangent Plane}

\begin{propositionbox}{Definition (Regularity, Unit Normal, Tangent Plane)}
The surface is \emph{regular} at $(u_0,v_0)$ if
$S_u \times S_v \neq \mathbf{0}$ there.

At a regular point, the \emph{unit normal vector} is
\[
  N(u_0,v_0)
  = \frac{S_u(u_0,v_0)\;\times\; S_v(u_0,v_0)}
         {\|S_u(u_0,v_0)\;\times\; S_v(u_0,v_0)\|}\,.
\]

The \emph{tangent plane} at $S(u_0,v_0)$ is the plane through
$S(u_0,v_0)$ spanned by $S_u$ and $S_v$:
\[
  \Pi\;=\;\bigl\{\,S(u_0,v_0) + \alpha\,S_u + \beta\,S_v
               \;\bigm|\; \alpha,\beta\in\mathbb{R}\,\bigr\}.
\]
Its equation in terms of position $\mathbf{p}$ is
\[
  N\cdot\bigl(\mathbf{p} - S(u_0,v_0)\bigr) = 0.
\]
\end{propositionbox}

\begin{remark*}[Geometric meaning of $S_u \times S_v$]
The magnitude $\|S_u \times S_v\|$ equals the area of the
parallelogram spanned by $S_u$ and $S_v$, which is the
\emph{local area element} of the surface.
Where $\|S_u \times S_v\| = 0$ the surface is degenerate
(the two tangent directions are parallel or one vanishes),
and the normal is not defined.
\end{remark*}

% ---------------------------------------------------------
% Figure
% ---------------------------------------------------------
\vspace{1em}
\begin{center}
\begin{tikzpicture}[
  >=Stealth,
  dot/.style={circle, fill, inner sep=1.8pt},
  lbl/.style={font=\small},
  slbl/.style={font=\scriptsize},
  vec/.style={->, line width=1.1pt},
  thinvec/.style={->, line width=0.7pt, gray!60},
]

% ---------------------------------------------------------------
% Coordinate helpers
% We draw a surface patch as a shaded quadrilateral in 3D
% using a mild oblique projection.
% Projection: (x,y,z) -> (x + 0.35*y, z + 0.20*y)
% ---------------------------------------------------------------

% Corner control points of the surface patch (in projection)
% uv-domain [0,1]x[0,1]; corners map to 3D:
%   (0,0)->P00, (1,0)->P10, (0,1)->P01, (1,1)->P11

\coordinate (P00) at (0.00, 0.00);   % u=0,v=0  (front-left)
\coordinate (P10) at (3.20, 0.30);   % u=1,v=0  (front-right)
\coordinate (P01) at (0.90, 2.60);   % u=0,v=1  (back-left)
\coordinate (P11) at (4.10, 2.85);   % u=1,v=1  (back-right)

% Surface point S(u0,v0) --- interior evaluation point
% roughly u0=0.45, v0=0.50
\coordinate (Suv) at (2.05, 1.38);

% -- Shaded surface patch ------------------------------------
\fill[blue!8]
  (P00) .. controls (1.0, -0.20) and (2.2, 0.10) .. (P10)
        .. controls (3.8,  1.20) and (4.3, 1.90) .. (P11)
        .. controls (2.7,  2.70) and (1.6, 2.65) .. (P01)
        .. controls (0.3,  1.50) and (-0.1,0.80) .. cycle;

\draw[blue!45!black, line width=1.0pt]
  (P00) .. controls (1.0, -0.20) and (2.2, 0.10) .. (P10)
        .. controls (3.8,  1.20) and (4.3, 1.90) .. (P11)
        .. controls (2.7,  2.70) and (1.6, 2.65) .. (P01)
        .. controls (0.3,  1.50) and (-0.1,0.80) .. cycle;

% -- Isoparametric u-curve through S(u0,v0) ------------------
% v = v0 $\approx$ 0.50: curve from left edge to right edge
\draw[red!60!black, line width=1.0pt, dashed]
  (0.45, 1.30) .. controls (1.2, 1.10) and (1.6, 1.30) ..
  (Suv) .. controls (2.5, 1.45) and (3.2, 1.50) .. (3.65, 1.60);
\node[slbl, text=red!60!black, anchor=north] at (3.65, 1.52)
  {$u$-curve};

% -- Isoparametric v-curve through S(u0,v0) ------------------
% u = u0 $\approx$ 0.45: curve from bottom edge to top edge
\draw[green!45!black, line width=1.0pt, dashed]
  (1.85, 0.14) .. controls (1.90, 0.65) and (2.00, 1.00) ..
  (Suv) .. controls (2.10, 1.75) and (2.20, 2.25) .. (2.30, 2.72);
\node[slbl, text=green!40!black, anchor=west] at (2.32, 2.72)
  {$v$-curve};

% -- Surface point S(u0,v0) -----------------------------------
\node[dot, fill=black] at (Suv) {};
\node[lbl, anchor=north east] at ($(Suv) + (-0.10, -0.08)$)
  {$S(u_0,v_0)$};

% -- Tangent vector S_u (red, along u-curve direction) --------
\coordinate (Su) at ($(Suv) + (1.55, 0.18)$);
\draw[vec, red!65!black] (Suv) -- (Su);
\node[lbl, text=red!65!black, anchor=west] at ($(Su) + (0.05, 0.0)$)
  {$S_u$};

% -- Tangent vector S_v (green, along v-curve direction) ------
\coordinate (Sv) at ($(Suv) + (0.24, 1.45)$);
\draw[vec, green!50!black] (Suv) -- (Sv);
\node[lbl, text=green!45!black, anchor=south] at ($(Sv) + (0.08, 0.05)$)
  {$S_v$};

% -- Normal vector N = S_u x S_v (blue, out of surface) -------
% Normal points up and slightly left out of the tangent plane
\coordinate (Nvec) at ($(Suv) + (-0.55, 2.10)$);
\draw[vec, blue!65!black, line width=1.4pt] (Suv) -- (Nvec);
\node[lbl, text=blue!65!black, anchor=south east]
  at ($(Nvec) + (-0.05, 0.05)$) {$N = \dfrac{S_u \times S_v}{\|S_u \times S_v\|}$};

% -- Tangent plane (shaded parallelogram) ---------------------
% Span S_u and S_v at the base point
\coordinate (TPa) at ($(Suv) + (1.55, 0.18) + (0.24, 1.45)$);  % Su + Sv
\fill[orange!10, opacity=0.70]
  (Suv) -- (Su) -- (TPa) -- (Sv) -- cycle;
\draw[orange!50!black, line width=0.6pt, dashed]
  (Suv) -- (Su) -- (TPa) -- (Sv) -- cycle;
\node[slbl, text=orange!55!black] at ($(Suv)!0.5!(TPa) + (0.15, -0.15)$)
  {tangent plane};

% -- Parameter domain box (lower-left inset) ------------------
\begin{scope}[shift={(-0.3, -2.8)}]
  % uv box
  \draw[gray!50, line width=0.8pt]
    (0,0) rectangle (1.8, 1.4);
  \node[slbl, text=gray!50!black] at (0.9, -0.22) {parameter domain};

  % axes labels
  \node[slbl] at (1.95, 0)    {$u$};
  \node[slbl] at (-0.12, 1.4) {$v$};
  \draw[thinvec] (0,0) -- (2.05,0);
  \draw[thinvec] (0,0) -- (0,1.55);

  % mark (u0,v0) in domain
  \node[dot, fill=black, inner sep=1.3pt] at (0.81, 0.70) {};
  \node[slbl, anchor=north west] at (0.84, 0.68) {$(u_0,v_0)$};

  % tick marks at 0 and 1
  \draw (0, 0.03) -- (0,-0.06) node[slbl, below] {$0$};
  \draw (1.8, 0.03) -- (1.8,-0.06) node[slbl, below] {$1$};
  \draw (0.03, 0) -- (-0.06, 0) node[slbl, left] {$0$};
  \draw (0.03,1.4) -- (-0.06,1.4) node[slbl, left] {$1$};
\end{scope}

% -- Arrow: domain -> surface ---------------------------------
\draw[->, gray!55, line width=0.9pt]
  (1.25, -1.90) .. controls (1.50, -0.90) and (1.60, -0.20) .. (Suv);
\node[slbl, text=gray!55!black] at (1.85, -1.25) {$S(u,v)$};

% -- Surface label ---------------------------------------------
\node[lbl, text=blue!55!black, font=\small\itshape]
  at (3.70, 2.80) {$S$};

\end{tikzpicture}

\end{center}

\begin{center}
\small\itshape
Figure: A parametric surface patch $S(u,v)$ with its parameter domain (lower left).
At the point $S(u_0,v_0)$, the partial derivative vectors $S_u$ (red) and $S_v$ (green)
are tangent to the isoparametric curves; together they span the tangent plane (orange).
The unit normal $N = S_u \times S_v\,/\,\|S_u \times S_v\|$ (blue) is perpendicular
to the tangent plane.
\end{center}

% ---------------------------------------------------------
% Status
% ---------------------------------------------------------
\vspace{1em}
\begin{tcolorbox}[
  colback=gray!6, colframe=gray!40, arc=2pt,
  left=8pt, right=8pt, top=6pt, bottom=6pt,
  title={\small\textbf{Status: Active}},
  fonttitle=\small\bfseries
]
\noindent This section is actively expanding alongside work in Piegl \& Tiller,
\textit{The NURBS Book}.
Subsequent sections on B\'{e}zier curves, B-splines, NURBS, and continuity
conditions follow.
\end{tcolorbox}
